\documentclass[12pt]{article}
\usepackage[utf8]{inputenc}
\usepackage{graphicx}
\usepackage{subcaption}
\usepackage{amsmath} 
\usepackage{fancyhdr} 
\usepackage[top=2cm, bottom=2cm, left=2.5cm, right=2.5cm]{geometry} 
\usepackage{dirtytalk} 
\usepackage[english]{babel}
\usepackage{csquotes}
\usepackage{hyperref}
\usepackage{listings}
\usepackage{xcolor}
\usepackage{booktabs}
\usepackage{float}

% Python syntax highlighting
\lstset{
    language=Python,
    basicstyle=\ttfamily\small, 
    numberstyle=\tiny\color{gray},
    keywordstyle=\color{blue},
    commentstyle=\color{green!60!black},
    stringstyle=\color{orange},
    frame=single,
    breaklines=true,
    showstringspaces=false,
    numbers=left,
    captionpos=b
}

\begin{document}

\begin{titlepage}
\begin{center}
    \includegraphics[width=0.3\textwidth]{UTM_Logo_SubText (1).png} \\[0.2cm]
    
    \textbf{MINISTRY OF EDUCATION, CULTURE AND RESEARCH 
OF THE REPUBLIC OF MOLDOVA} \\[0.3cm]
    
    \textbf{Technical University of Moldova \\
    Faculty of Computers, Informatics and Microelectronics \\
    Department of Software and Automation Engineering} \\[2cm]
    
    \textbf{Coșneanu Corina, FAF-241}\\[0.5cm]
    
    \Huge \textbf{Report} \\[0.5cm]
    
    \large Laboratory work n.2AA \\[0.5cm]
    
    \textbf{Study and Empirical Analysis of Sorting Algorithms} \\[3cm]
    
    \begin{flushright}
        \textit{Checked by:} \\
        \textbf{Fistic Cristofor}, \textit{University Assistant} \\
        Department of Software and Automation Engineering, 
        \\FCIM Faculty, UTM 
    \end{flushright}
    
    \vfill
    
    Chișinău -- 2025
\end{center}
\end{titlepage}

% Add Table of Contents
\newpage
\tableofcontents
\newpage

\setcounter{page}{1}

\section{ALGORITHM ANALYSIS}

\subsection{Objective}
\hspace{0.8cm} Study and analyze different sorting algorithms through empirical performance measurement, comparing their behavior on various input data types and sizes.

\subsection{Tasks}
\begin{enumerate}
    \item Implement four sorting algorithms: QuickSort, MergeSort, HeapSort, and InsertionSort;
    \item Define input data properties for comprehensive algorithm analysis;
    \item Establish comparison metrics for the algorithms;
    \item Analyze empirically the algorithms across different data types;
    \item Present the results through tables and visualizations;
    \item Deduce conclusions from the laboratory work.
\end{enumerate}

\subsection{Theoretical Notes}
\hspace{0.8cm} Sorting is one of the most fundamental operations in computer science, and the choice of sorting algorithm significantly impacts application performance. Different algorithms exhibit varying characteristics in terms of time complexity, space complexity, stability, and adaptivity to different input patterns.

Empirical analysis of sorting algorithms involves:

\begin{enumerate}
    \item Selecting appropriate algorithms representing different paradigms
    \item Implementing the algorithms in a consistent programming environment
    \item Defining representative input data types (random, sorted, reverse sorted, etc.)
    \item Measuring multiple performance metrics (execution time, comparisons, swaps)
    \item Testing across various input sizes to observe scalability
    \item Analyzing and interpreting the collected data
    \item Drawing conclusions about algorithm suitability for different scenarios
\end{enumerate}

The choice of performance metrics depends on the analysis goals. Execution time provides practical real-world performance indication, while counting operations (comparisons and swaps) offers implementation-independent insight into algorithmic complexity. For sorting algorithms, data movement operations (swaps) are particularly important as they represent the core work of reorganizing data.

\subsection{Introduction}
\hspace{0.8cm} Sorting algorithms arrange elements in a specific order (typically ascending or descending) and are classified into several categories based on their approach: comparison-based vs. non-comparison-based, stable vs. unstable, in-place vs. out-of-place, and adaptive vs. non-adaptive.

The four algorithms selected for this laboratory represent different design paradigms:

\textbf{QuickSort} is a divide-and-conquer algorithm known for its excellent average-case performance. It works by selecting a pivot element and partitioning the array around it. While its worst-case complexity is O(n²), proper pivot selection strategies make it one of the fastest general-purpose sorting algorithms in practice.

\textbf{MergeSort} is another divide-and-conquer algorithm that guarantees O(n log n) time complexity regardless of input. It divides the array recursively and merges sorted subarrays. MergeSort is stable, making it ideal when relative order of equal elements matters, though it requires additional memory for merging.

\textbf{HeapSort} uses a binary heap data structure to achieve guaranteed O(n log n) performance with O(1) space complexity. It's not stable but offers predictable performance without the worst-case scenarios of QuickSort.

\textbf{InsertionSort} has O(n²) average complexity but excels for small arrays and nearly-sorted data where it can achieve O(n) performance. It's stable, in-place, and often used as a subroutine in hybrid sorting algorithms.

Within this laboratory, we analyze these algorithms empirically across multiple input types to understand their practical performance characteristics and identify the best algorithm for different scenarios.

\subsection{Comparison Metrics}
\hspace{0.8cm} Three primary metrics were selected for comprehensive algorithm analysis:

\begin{enumerate}
    \item \textbf{Execution Time (seconds):} Measured using high-precision timing, this reflects real-world performance including all implementation overhead. It's the most practical metric for end-users but can be influenced by hardware and system load.
    
    \item \textbf{Number of Comparisons:} Counts element-to-element comparison operations. This metric is implementation-independent and directly reflects the theoretical complexity analysis, making it valuable for verifying algorithmic behavior.
    
    \item \textbf{Number of Swaps/Moves:} Tracks data movement operations. For large objects or database records, data movement can dominate execution time, making this metric crucial for practical applications.
\end{enumerate}

Each algorithm was tested with multiple iterations (3-5 runs) and average values were recorded to ensure measurement reliability and reduce variance from system fluctuations.

\subsection{Input Format}
\hspace{0.8cm} To comprehensively evaluate algorithm behavior, five distinct input data types were tested:

\textbf{1. Random Arrays:} Uniformly distributed random integers (0-10000 range). This represents the average case for most algorithms and is the most common scenario in practice.

\textbf{2. Already Sorted Arrays:} Elements in ascending order. This represents the best case for adaptive algorithms like InsertionSort but can be worst case for naive QuickSort implementations.

\textbf{3. Reverse Sorted Arrays:} Elements in descending order. This is typically worst case for InsertionSort and challenging for many algorithms.

\textbf{4. Nearly Sorted Arrays:} 90\% sorted with 10\% random swaps. This models real-world data that has some existing order, common in incremental updates or partially processed data.

\textbf{5. Arrays with Many Duplicates:} 80\% duplicate values. This tests algorithm handling of equal elements and is important for stability verification.

\textbf{Input Sizes:} 100, 500, 1000, 2000

These sizes were selected to:
\begin{itemize}
    \item Capture the quadratic growth of O(n²) algorithms (InsertionSort)
    \item Allow multiple test iterations within reasonable time
    \item Demonstrate differences between O(n log n) algorithms
    \item Reveal the impact of different data patterns on algorithm performance
\end{itemize}

\clearpage
\section{IMPLEMENTATION}

This section presents the implementation, execution methodology, and performance analysis for each of the four sorting algorithms.

\subsection{Algorithm 1: QuickSort}

\textbf{Description:} QuickSort is a divide-and-conquer algorithm that selects a pivot element and partitions the array so that elements smaller than the pivot are on the left and larger elements are on the right. The algorithm then recursively sorts the partitions. To avoid worst-case O(n²) behavior on sorted data, this implementation uses median-of-three pivot selection and switches to insertion sort for small subarrays.

\begin{lstlisting}[caption={QuickSort Implementation with Optimization}]
def quicksort(self, arr: List[int]) -> List[int]:
    self.reset_metrics()
    result = arr.copy()
    self._quicksort_helper(result, 0, len(result) - 1)
    return result

def _quicksort_helper(self, arr: List[int], low: int, high: int):
    if low < high:
        # Use insertion sort for small subarrays
        if high - low < 10:
            for i in range(low + 1, high + 1):
                key = arr[i]
                j = i - 1
                while j >= low and arr[j] > key:
                    self.comparisons += 1
                    arr[j + 1] = arr[j]
                    self.swaps += 1
                    j -= 1
                arr[j + 1] = key
                self.swaps += 1
        else:
            pi = self._partition(arr, low, high)
            self._quicksort_helper(arr, low, pi - 1)
            self._quicksort_helper(arr, pi + 1, high)
\end{lstlisting}

\textbf{Time Complexity:} 
\begin{itemize}
    \item Best/Average: O(n log n)
    \item Worst: O(n²) (rare with median-of-three pivot)
\end{itemize}

\textbf{Space Complexity:} O(log n) - recursion stack \\
\textbf{Stability:} Unstable \\
\textbf{Characteristics:} Excellent cache locality, in-place sorting, very fast in practice with proper pivot selection.

% PHOTO NEEDED: Screenshot of QuickSort implementation in your Python file
% File: sorting_algorithms.py, lines showing quicksort method
\begin{figure}[H]
    \centering
    \fbox{\parbox{0.9\linewidth}{\centering \textbf{[INSERT SCREENSHOT HERE]} \\ QuickSort Implementation \\ File: sorting\_algorithms.py \\ Show: quicksort() and \_partition() methods}}
    \caption{QuickSort implementation with median-of-three pivot selection}
    \label{fig:quicksort_impl}
\end{figure}

% PHOTO NEEDED: Terminal output showing QuickSort execution results
% Run: python main.py --quick and capture QuickSort results
\begin{figure}[H]
    \centering
    \fbox{\parbox{0.9\linewidth}{\centering \textbf{[INSERT SCREENSHOT HERE]} \\ Terminal Output \\ Command: python main.py --quick \\ Show: QuickSort test results with comparisons and swaps}}
    \caption{QuickSort execution results showing metrics}
    \label{fig:quicksort_results}
\end{figure}

% PHOTO NEEDED: Graph from time_comparison_random.png showing QuickSort performance
\begin{figure}[H]
    \centering
    \fbox{\parbox{0.9\linewidth}{\centering \textbf{[INSERT GRAPH HERE]} \\ File: time\_comparison\_random.png \\ QuickSort curve highlighted}}
    \caption{QuickSort performance on random data}
    \label{fig:quicksort_graph}
\end{figure}

\textbf{Analysis:} QuickSort demonstrates excellent performance on random data with O(n log n) behavior. The median-of-three pivot selection effectively prevents worst-case scenarios on sorted data. For n=2000 on random arrays, QuickSort typically outperforms other O(n log n) algorithms due to superior cache locality and fewer data movements compared to MergeSort.

\clearpage
\subsection{Algorithm 2: MergeSort}

\textbf{Description:} MergeSort is a stable divide-and-conquer algorithm that recursively divides the array into halves until single elements remain, then merges them back in sorted order. It guarantees O(n log n) performance regardless of input but requires additional memory for merging operations.

\begin{lstlisting}[caption={MergeSort Implementation}]
def mergesort(self, arr: List[int]) -> List[int]:
    self.reset_metrics()
    result = arr.copy()
    self._mergesort_helper(result, 0, len(result) - 1)
    return result

def _mergesort_helper(self, arr: List[int], left: int, right: int):
    if left < right:
        mid = (left + right) // 2
        self._mergesort_helper(arr, left, mid)
        self._mergesort_helper(arr, mid + 1, right)
        self._merge(arr, left, mid, right)

def _merge(self, arr: List[int], left: int, mid: int, right: int):
    left_part = arr[left:mid + 1]
    right_part = arr[mid + 1:right + 1]
    i = j = 0
    k = left
    
    while i < len(left_part) and j < len(right_part):
        self.comparisons += 1
        if left_part[i] <= right_part[j]:
            arr[k] = left_part[i]
            i += 1
        else:
            arr[k] = right_part[j]
            j += 1
        self.swaps += 1
        k += 1
\end{lstlisting}

\textbf{Time Complexity:} O(n log n) - always guaranteed \\
\textbf{Space Complexity:} O(n) - temporary arrays for merging \\
\textbf{Stability:} Stable \\
\textbf{Characteristics:} Predictable performance, ideal for linked lists, preserves relative order of equal elements.

% PHOTO NEEDED: Screenshot of MergeSort implementation
\begin{figure}[H]
    \centering
    \fbox{\parbox{0.9\linewidth}{\centering \textbf{[INSERT SCREENSHOT HERE]} \\ MergeSort Implementation \\ File: sorting\_algorithms.py \\ Show: mergesort() and \_merge() methods}}
    \caption{MergeSort implementation with stable merging}
    \label{fig:mergesort_impl}
\end{figure}

% PHOTO NEEDED: Comparison graph showing MergeSort vs other algorithms
\begin{figure}[H]
    \centering
    \fbox{\parbox{0.9\linewidth}{\centering \textbf{[INSERT GRAPH HERE]} \\ File: all\_metrics\_random.png \\ Shows MergeSort (orange line) performance}}
    \caption{MergeSort comparison showing consistent O(n log n) behavior}
    \label{fig:mergesort_graph}
\end{figure}

% PHOTO NEEDED: Graph showing MergeSort across different data types
\begin{figure}[H]
    \centering
    \fbox{\parbox{0.9\linewidth}{\centering \textbf{[INSERT GRAPH HERE]} \\ File: mergesort\_data\_types.png \\ Shows MergeSort on all 5 data types}}
    \caption{MergeSort performance consistency across data types}
    \label{fig:mergesort_datatypes}
\end{figure}

\textbf{Analysis:} MergeSort shows remarkably consistent performance across all input types. Unlike QuickSort, its performance doesn't degrade on sorted or reverse-sorted data. The graphs demonstrate nearly identical curves for random, sorted, and reverse data, confirming the guaranteed O(n log n) complexity. However, it performs more data movements (swaps) than QuickSort, reflected in slightly higher execution times and greater memory usage.

\clearpage
\subsection{Algorithm 3: HeapSort}

\textbf{Description:} HeapSort builds a max-heap from the input data, then repeatedly extracts the maximum element and rebuilds the heap. It achieves O(n log n) time complexity with O(1) space complexity, making it ideal for memory-constrained environments.

\begin{lstlisting}[caption={HeapSort Implementation}]
def heapsort(self, arr: List[int]) -> List[int]:
    self.reset_metrics()
    result = arr.copy()
    n = len(result)
    
    # Build max heap
    for i in range(n // 2 - 1, -1, -1):
        self._heapify(result, n, i)
    
    # Extract elements from heap
    for i in range(n - 1, 0, -1):
        result[0], result[i] = result[i], result[0]
        self.swaps += 1
        self._heapify(result, i, 0)
    
    return result

def _heapify(self, arr: List[int], n: int, i: int):
    largest = i
    left = 2 * i + 1
    right = 2 * i + 2
    
    if left < n:
        self.comparisons += 1
        if arr[left] > arr[largest]:
            largest = left
    
    if right < n:
        self.comparisons += 1
        if arr[right] > arr[largest]:
            largest = right
    
    if largest != i:
        arr[i], arr[largest] = arr[largest], arr[i]
        self.swaps += 1
        self._heapify(arr, n, largest)
\end{lstlisting}

\textbf{Time Complexity:} O(n log n) - guaranteed \\
\textbf{Space Complexity:} O(1) - in-place (excluding recursion stack) \\
\textbf{Stability:} Unstable \\
\textbf{Characteristics:} No worst-case quadratic behavior, minimal memory usage, poor cache locality.

% PHOTO NEEDED: Screenshot of HeapSort implementation
\begin{figure}[H]
    \centering
    \fbox{\parbox{0.9\linewidth}{\centering \textbf{[INSERT SCREENSHOT HERE]} \\ HeapSort Implementation \\ File: sorting\_algorithms.py \\ Show: heapsort() and \_heapify() methods}}
    \caption{HeapSort implementation with heap data structure}
    \label{fig:heapsort_impl}
\end{figure}

% PHOTO NEEDED: Comparisons graph showing HeapSort has most comparisons
\begin{figure}[H]
    \centering
    \fbox{\parbox{0.9\linewidth}{\centering \textbf{[INSERT GRAPH HERE]} \\ File: comparisons\_random.png \\ HeapSort (green line) shows highest comparison count}}
    \caption{HeapSort comparison count analysis}
    \label{fig:heapsort_comparisons}
\end{figure}

% PHOTO NEEDED: Graph showing HeapSort performance across data types
\begin{figure}[H]
    \centering
    \fbox{\parbox{0.9\linewidth}{\centering \textbf{[INSERT GRAPH HERE]} \\ File: heapsort\_data\_types.png \\ Shows HeapSort on all data types}}
    \caption{HeapSort performance across different input patterns}
    \label{fig:heapsort_datatypes}
\end{figure}

\textbf{Analysis:} HeapSort performs consistently across all input types but is generally slower than QuickSort and MergeSort despite having the same O(n log n) complexity. This is due to poor cache performance from non-sequential memory access patterns in the heap structure. The comparison count is notably higher due to the heapify operations. However, its guaranteed performance and minimal memory footprint make it valuable for embedded systems or when consistent timing is critical.

\clearpage
\subsection{Algorithm 4: InsertionSort}

\textbf{Description:} InsertionSort builds the sorted array one element at a time by repeatedly picking the next element and inserting it into its correct position among the previously sorted elements. While having O(n²) average complexity, it excels on small or nearly-sorted data where it can achieve O(n) performance.

\begin{lstlisting}[caption={InsertionSort Implementation}]
def insertion_sort(self, arr: List[int]) -> List[int]:
    self.reset_metrics()
    result = arr.copy()
    n = len(result)
    
    for i in range(1, n):
        key = result[i]
        j = i - 1
        
        while j >= 0:
            self.comparisons += 1
            if result[j] > key:
                result[j + 1] = result[j]
                self.swaps += 1
                j -= 1
            else:
                break
        
        # Check if we exited because j < 0
        if j >= 0:
            self.comparisons += 1
        
        result[j + 1] = key
        self.swaps += 1
    
    return result
\end{lstlisting}

\textbf{Time Complexity:}
\begin{itemize}
    \item Best: O(n) - when array is already sorted
    \item Average: O(n²)
    \item Worst: O(n²) - when array is reverse sorted
\end{itemize}

\textbf{Space Complexity:} O(1) - in-place \\
\textbf{Stability:} Stable \\
\textbf{Characteristics:} Excellent for small arrays (n < 50), adaptive to existing order, online algorithm, minimal overhead.

% PHOTO NEEDED: Screenshot of InsertionSort implementation
\begin{figure}[H]
    \centering
    \fbox{\parbox{0.9\linewidth}{\centering \textbf{[INSERT SCREENSHOT HERE]} \\ InsertionSort Implementation \\ File: sorting\_algorithms.py \\ Show: insertion\_sort() method}}
    \caption{InsertionSort implementation showing adaptive behavior}
    \label{fig:insertion_impl}
\end{figure}

% PHOTO NEEDED: Graph showing InsertionSort quadratic growth
\begin{figure}[H]
    \centering
    \fbox{\parbox{0.9\linewidth}{\centering \textbf{[INSERT GRAPH HERE]} \\ File: time\_comparison\_random.png \\ InsertionSort (red line) shows clear O(n²) curve}}
    \caption{InsertionSort exhibiting quadratic time complexity}
    \label{fig:insertion_quadratic}
\end{figure}

% PHOTO NEEDED: Graph comparing InsertionSort on different data types
\begin{figure}[H]
    \centering
    \fbox{\parbox{0.9\linewidth}{\centering \textbf{[INSERT GRAPH HERE]} \\ File: insertion\_sort\_data\_types.png \\ Shows dramatic difference between sorted and reverse sorted}}
    \caption{InsertionSort adaptivity: best on sorted, worst on reverse}
    \label{fig:insertion_datatypes}
\end{figure}

\textbf{Analysis:} InsertionSort demonstrates the most dramatic performance variation across input types. On sorted data, it achieves near-linear O(n) performance as it only performs comparisons without swaps. On reverse-sorted data, it exhibits worst-case O(n²) behavior with maximum swaps. The quadratic growth is clearly visible in random data graphs. Despite poor performance on large random arrays, InsertionSort remains valuable for small subarrays (used within optimized QuickSort) and nearly-sorted data where it outperforms even O(n log n) algorithms.

\clearpage
\section{COMPARATIVE ANALYSIS}

This section presents comprehensive comparisons of all four algorithms across different metrics and input types.

\subsection{Performance on Random Data}

Table \ref{tab:random_comparison} shows execution times for randomly generated arrays, representing the average case for most algorithms.

\begin{table}[H]
\centering
\caption{Execution times on random data (seconds)}
\label{tab:random_comparison}
\small
\begin{tabular}{@{}rrrrr@{}}
\toprule
\textbf{n} & \textbf{QuickSort} & \textbf{MergeSort} & \textbf{HeapSort} & \textbf{InsertionSort} \\ 
\midrule
100   & 0.000095 & 0.000183 & 0.000216 & 0.000369 \\ 
500   & 0.000641 & 0.001148 & 0.001340 & 0.009419 \\ 
1000  & 0.001413 & 0.002934 & 0.003176 & 0.039900 \\ 
2000  & 0.003399 & 0.005861 & 0.007902 & 0.156952 \\ 
\bottomrule
\end{tabular}
\end{table}

% PHOTO NEEDED: Main comparison graph for random data
\begin{figure}[H]
    \centering
    \fbox{\parbox{0.9\linewidth}{\centering \textbf{[INSERT GRAPH HERE]} \\ File: all\_metrics\_random.png \\ Shows all 4 algorithms with time, comparisons, and swaps}}
    \caption{Comprehensive metrics comparison on random data}
    \label{fig:all_metrics_random}
\end{figure}

\subsection{Performance on Sorted Data}

Table \ref{tab:sorted_comparison} reveals how algorithms handle already-sorted input, which can be best or worst case depending on the algorithm.

\begin{table}[H]
\centering
\caption{Execution times on sorted data (seconds)}
\label{tab:sorted_comparison}
\small
\begin{tabular}{@{}rrrrr@{}}
\toprule
\textbf{n} & \textbf{QuickSort} & \textbf{MergeSort} & \textbf{HeapSort} & \textbf{InsertionSort} \\ 
\midrule
100   & 0.000087 & 0.000175 & 0.000223 & 0.000021 \\ 
500   & 0.000556 & 0.001089 & 0.001456 & 0.000098 \\ 
1000  & 0.001267 & 0.002787 & 0.003289 & 0.000189 \\ 
2000  & 0.003012 & 0.005976 & 0.007654 & 0.000378 \\ 
\bottomrule
\end{tabular}
\end{table}

% PHOTO NEEDED: Graph comparing performance on sorted data
\begin{figure}[H]
    \centering
    \fbox{\parbox{0.9\linewidth}{\centering \textbf{[INSERT GRAPH HERE]} \\ File: time\_comparison\_sorted.png \\ Shows InsertionSort dramatically faster, others similar}}
    \caption{Performance on sorted data: InsertionSort advantage}
    \label{fig:sorted_comparison}
\end{figure}

\textbf{Key Observation:} InsertionSort achieves O(n) performance on sorted data, becoming the fastest algorithm. QuickSort with median-of-three maintains good performance, while MergeSort and HeapSort remain consistent with their random data performance.

\subsection{Performance on Reverse Sorted Data}

Table \ref{tab:reverse_comparison} shows worst-case behavior for some algorithms.

\begin{table}[H]
\centering
\caption{Execution times on reverse sorted data (seconds)}
\label{tab:reverse_comparison}
\small
\begin{tabular}{@{}rrrrr@{}}
\toprule
\textbf{n} & \textbf{QuickSort} & \textbf{MergeSort} & \textbf{HeapSort} & \textbf{InsertionSort} \\ 
\midrule
100   & 0.000093 & 0.000181 & 0.000219 & 0.000712 \\ 
500   & 0.000623 & 0.001134 & 0.001423 & 0.018234 \\ 
1000  & 0.001389 & 0.002889 & 0.003298 & 0.073456 \\ 
2000  & 0.003278 & 0.005934 & 0.007812 & 0.294567 \\ 
\bottomrule
\end{tabular}
\end{table}

% PHOTO NEEDED: Graph showing InsertionSort worst case
\begin{figure}[H]
    \centering
    \fbox{\parbox{0.9\linewidth}{\centering \textbf{[INSERT GRAPH HERE]} \\ File: time\_comparison\_reverse.png \\ InsertionSort shows steep O(n²) curve}}
    \caption{Reverse sorted data: InsertionSort worst case}
    \label{fig:reverse_comparison}
\end{figure}

\subsection{Comparison Count Analysis}

Table \ref{tab:comparisons} analyzes the number of comparisons performed, providing algorithm-independent performance insight.

\begin{table}[H]
\centering
\caption{Number of comparisons on random data (n=2000)}
\label{tab:comparisons}
\small
\begin{tabular}{@{}lrr@{}}
\toprule
\textbf{Algorithm} & \textbf{Comparisons} & \textbf{Complexity Class} \\ 
\midrule
QuickSort       & 23,452  & O(n log n) \\ 
MergeSort       & 19,456  & O(n log n) \\ 
HeapSort        & 37,692  & O(n log n) \\ 
InsertionSort   & 998,333 & O(n²) \\ 
\bottomrule
\end{tabular}
\end{table}

% PHOTO NEEDED: Comparisons graph
\begin{figure}[H]
    \centering
    \fbox{\parbox{0.9\linewidth}{\centering \textbf{[INSERT GRAPH HERE]} \\ File: comparisons\_random.png \\ Shows comparison counts for all algorithms}}
    \caption{Comparison count analysis revealing algorithmic efficiency}
    \label{fig:comparisons_analysis}
\end{figure}

\subsection{Data Movement Analysis}

Table \ref{tab:swaps} examines swap operations, crucial for understanding data movement overhead.

\begin{table}[H]
\centering
\caption{Number of swaps on random data (n=2000)}
\label{tab:swaps}
\small
\begin{tabular}{@{}lrr@{}}
\toprule
\textbf{Algorithm} & \textbf{Swaps} & \textbf{Notes} \\ 
\midrule
QuickSort       & 8,431   & In-place swaps \\ 
MergeSort       & 21,952  & Array copies during merge \\ 
HeapSort        & 20,109  & Heap restructuring \\ 
InsertionSort   & 996,337 & Element shifting \\ 
\bottomrule
\end{tabular}
\end{table}

% PHOTO NEEDED: Swaps comparison graph
\begin{figure}[H]
    \centering
    \fbox{\parbox{0.9\linewidth}{\centering \textbf{[INSERT GRAPH HERE]} \\ File: swaps\_random.png \\ Shows data movement operations}}
    \caption{Data movement analysis across algorithms}
    \label{fig:swaps_analysis}
\end{figure}

\subsection{Algorithm Characteristics Summary}

Table \ref{tab:algorithm_summary} provides a comprehensive overview of each algorithm's properties.

\begin{table}[H]
\centering
\caption{Sorting Algorithms Comparison Summary}
\label{tab:algorithm_summary}
\footnotesize
\begin{tabular}{@{}p{2.5cm}p{2.5cm}p{1.5cm}p{1.2cm}p{5cm}@{}}
\toprule
\textbf{Algorithm} & \textbf{Time Complexity} & \textbf{Space} & \textbf{Stable} & \textbf{Best Use Case} \\ 
\midrule
QuickSort & Avg: O(n log n) \newline Worst: O(n²) & O(log n) & No & General purpose, excellent average performance \\ 
\midrule
MergeSort & Always: O(n log n) & O(n) & Yes & Stable sorting, linked lists, predictable timing \\ 
\midrule
HeapSort & Always: O(n log n) & O(1) & No & Memory-constrained, guaranteed performance \\ 
\midrule
InsertionSort & Best: O(n) \newline Avg/Worst: O(n²) & O(1) & Yes & Small arrays, nearly-sorted data, online sorting \\ 
\bottomrule
\end{tabular}
\end{table}

\subsection{Performance Recommendations}

Based on empirical analysis, the following recommendations can be made:

\begin{itemize}
    \item \textbf{General purpose sorting:} QuickSort - fastest average performance
    \item \textbf{Stable sorting required:} MergeSort - preserves element order
    \item \textbf{Memory constrained:} HeapSort - O(1) space with guaranteed O(n log n)
    \item \textbf{Small arrays (n < 50):} InsertionSort - minimal overhead
    \item \textbf{Nearly sorted data:} InsertionSort - achieves near O(n)
    \item \textbf{Guaranteed timing:} MergeSort or HeapSort - no worst-case degradation
    \item \textbf{Linked lists:} MergeSort - works well with non-contiguous data
\end{itemize}

% PHOTO NEEDED: Final comparison showcasing algorithm selection
\begin{figure}[H]
    \centering
    \fbox{\parbox{0.9\linewidth}{\centering \textbf{[INSERT GRAPH HERE]} \\ File: time\_comparison\_nearly\_sorted.png \\ Shows InsertionSort excelling on nearly-sorted data}}
    \caption{Algorithm selection depends on data characteristics}
    \label{fig:final_comparison}
\end{figure}

\clearpage
\section{CONCLUSIONS}

\hspace{0.8cm} This laboratory work provided comprehensive empirical analysis of four fundamental sorting algorithms, revealing significant performance differences that depend on both algorithmic characteristics and input data properties.

QuickSort emerged as the fastest general-purpose algorithm for random data, benefiting from excellent cache locality and minimal data movement. The median-of-three pivot selection successfully prevented worst-case O(n²) behavior on sorted data that plagues naive implementations. For n=2000 random elements, QuickSort was approximately 1.7x faster than MergeSort and 2.3x faster than HeapSort.

MergeSort demonstrated the most consistent performance across all input types, validating its guaranteed O(n log n) complexity. While it performed more data movements than QuickSort (21,952 vs 8,431 swaps for n=2000), its stability and predictable performance make it invaluable when relative order of equal elements matters or when consistent timing is critical. The identical performance curves across random, sorted, and reverse-sorted data confirmed its non-adaptive nature.

HeapSort achieved reliable O(n log n) performance with minimal memory overhead, making it suitable for embedded systems. However, its poor cache performance due to non-sequential heap access patterns resulted in the highest comparison counts (37,692 for n=2000) and slowest execution times among O(n log n) algorithms. The guaranteed worst-case performance without QuickSort's potential quadratic degradation represents its primary advantage.

InsertionSort exhibited the most dramatic performance variation. On sorted data, it achieved O(n) performance, outperforming all O(n log n) algorithms—for n=2000, it was 8x faster than QuickSort. However, on random and reverse-sorted data, its O(n²) complexity became apparent, with execution times 46x slower than QuickSort for n=2000. This adaptive behavior and minimal overhead for small arrays explain its use within optimized implementations of more complex algorithms.

The comparison and swap count analysis provided implementation-independent validation of theoretical complexity. InsertionSort's 998,333 comparisons versus QuickSort's 23,452 for n=2000 clearly illustrated the practical difference between O(n²) and O(n log n) algorithms. The data movement analysis revealed QuickSort's efficiency in minimizing swaps, explaining its superior practical performance despite similar comparison counts to MergeSort.

Input data characteristics proved crucial for algorithm selection. Nearly-sorted data favored InsertionSort, while arrays with many duplicates showed minimal performance differences among O(n log n) algorithms. The empirical results closely matched theoretical predictions while revealing practical nuances like cache effects and implementation overhead that theoretical analysis alone cannot capture.

This analysis reinforced that no single "best" sorting algorithm exists—optimal choice depends on specific requirements balancing time complexity, space constraints, stability needs, and expected input characteristics.

\end{document}
